% Preámbulo de las diapositivas
\usepackage{amsmath,amssymb}

% --- Símbolos de dificultad y de nivel -------------------------------------
% Latin Modern no tiene glifos para ○ ◐ ● ★ ◆, y XeTeX los descarta EN SILENCIO:
% el marcador de dificultad de cada problema desaparecía del PDF sin aviso.
% Se mapean a equivalentes matemáticos que sí existen.
\usepackage[nointegrals]{wasysym}
\usepackage{newunicodechar}
\newunicodechar{○}{\ensuremath{\Circle}}
\newunicodechar{◐}{\ensuremath{\LEFTcircle}}
\newunicodechar{●}{\ensuremath{\CIRCLE}}
% \bigstar y \blacklozenge NO valen aquí: la plantilla de pandoc carga
% unicode-math, que los reasigna a caracteres que la fuente matemática no tiene.
% Los dingbats de pifont son independientes de la fuente matemática.
\usepackage{pifont}
\newunicodechar{★}{\ding{72}}
\newunicodechar{◆}{\ding{117}}

% --- Fuente monoespaciada con cobertura Unicode ----------------------------
% Los diagramas ASCII y las plantillas del libro usan caracteres de dibujo de
% cajas (─ │ ├ └ ┬) y casillas (□). Latin Modern Mono no los tiene y XeTeX los
% descarta EN SILENCIO: los diagramas salían desmontados sin ningún aviso.
% DejaVu Sans Mono sí los cubre. Si no está instalada, se sigue sin ella.
\IfFontExistsTF{DejaVu Sans Mono}{%
  \setmonofont{DejaVu Sans Mono}[Scale=MatchLowercase]%
}{}

% Los mismos caracteres, cuando aparecen fuera de un bloque de código.
\newunicodechar{□}{\ding{113}}
\newunicodechar{✓}{\ding{51}}
\newunicodechar{⁻}{\ensuremath{^{-}}}
\newunicodechar{⁰}{\ensuremath{^{0}}}
\newunicodechar{⁴}{\ensuremath{^{4}}}
\newunicodechar{⁵}{\ensuremath{^{5}}}
\newunicodechar{⁶}{\ensuremath{^{6}}}
\newunicodechar{₀}{\ensuremath{_{0}}}
\newunicodechar{₂}{\ensuremath{_{2}}}
\newunicodechar{ᵢ}{\ensuremath{_{i}}}
\newunicodechar{ⱼ}{\ensuremath{_{j}}}
\newunicodechar{ṁ}{\ensuremath{\dot m}}

% Letras griegas y símbolos que aparecen en el texto corrido, fuera de modo
% matemático. Latin Modern Roman no las tiene y también se caían en silencio.
% Nada de upgreek: choca con unicode-math + microtype a través de mt-eur.cfg.
\newunicodechar{μ}{\ensuremath{\mu}}
\newunicodechar{σ}{\ensuremath{\sigma}}
\newunicodechar{π}{\ensuremath{\pi}}
\newunicodechar{ρ}{\ensuremath{\rho}}
\newunicodechar{α}{\ensuremath{\alpha}}
\newunicodechar{τ}{\ensuremath{\tau}}
\newunicodechar{χ}{\ensuremath{\chi}}
\newunicodechar{ε}{\ensuremath{\varepsilon}}
\newunicodechar{γ}{\ensuremath{\gamma}}
\newunicodechar{Δ}{\ensuremath{\Delta}}
\newunicodechar{Σ}{\ensuremath{\Sigma}}
\newunicodechar{∝}{\ensuremath{\propto}}
\newunicodechar{∞}{\ensuremath{\infty}}
\newunicodechar{≪}{\ensuremath{\ll}}
\newunicodechar{∼}{\ensuremath{\sim}}
\newunicodechar{⟨}{\ensuremath{\langle}}
\newunicodechar{⟩}{\ensuremath{\rangle}}

% Operadores que también aparecen sueltos en el texto corrido.
\newunicodechar{≈}{\ensuremath{\approx}}
\newunicodechar{≥}{\ensuremath{\geq}}
\newunicodechar{≤}{\ensuremath{\leq}}
\newunicodechar{÷}{\ensuremath{\div}}
\newunicodechar{−}{\ensuremath{-}}
\newunicodechar{√}{\ensuremath{\surd}}
\newunicodechar{≠}{\ensuremath{\neq}}

% Los caracteres de dibujo de cajas sólo aparecen dentro de bloques literales,
% donde los cubre DejaVu Sans Mono. Por si alguna vez aparecen en texto corrido,
% se fuerza la fuente monoespaciada. \symbol no vuelve a pasar por newunicodechar,
% así que no hay recursión.
\newunicodechar{─}{{\ttfamily\symbol{"2500}}}
\newunicodechar{│}{{\ttfamily\symbol{"2502}}}
\newunicodechar{┌}{{\ttfamily\symbol{"250C}}}
\newunicodechar{┐}{{\ttfamily\symbol{"2510}}}
\newunicodechar{└}{{\ttfamily\symbol{"2514}}}
\newunicodechar{┘}{{\ttfamily\symbol{"2518}}}
\newunicodechar{├}{{\ttfamily\symbol{"251C}}}
\newunicodechar{┤}{{\ttfamily\symbol{"2524}}}
\newunicodechar{┬}{{\ttfamily\symbol{"252C}}}
\newunicodechar{┴}{{\ttfamily\symbol{"2534}}}
\newunicodechar{┼}{{\ttfamily\symbol{"253C}}}
\definecolor{tinta}{HTML}{1B2A41}
\definecolor{azul}{HTML}{2F6FA8}
\definecolor{rojo}{HTML}{C1443C}
\definecolor{verde}{HTML}{3F8F6B}
\definecolor{ocre}{HTML}{D99A2B}
\definecolor{grisc}{HTML}{8A94A6}
\definecolor{papel}{HTML}{FBFAF7}
\setbeamercolor{normal text}{fg=tinta,bg=papel}
\setbeamercolor{alerted text}{fg=rojo}
\setbeamercolor{example text}{fg=verde}
\setbeamercolor{structure}{fg=azul}
\setbeamercolor{frametitle}{fg=papel,bg=tinta}
\setbeamercolor{title}{fg=tinta}
\setbeamercolor{progress bar}{fg=azul}
\setbeamerfont{frametitle}{size=\large,series=\bfseries}
\setbeamertemplate{navigation symbols}{}
\setbeamertemplate{itemize items}[circle]
\setbeamertemplate{footline}{%
  \hbox{\begin{beamercolorbox}[wd=\paperwidth,ht=2.4ex,dp=1.2ex,leftskip=.3cm,rightskip=.3cm]{}%
  \tiny\color{grisc}\insertshorttitle\hfill\insertframenumber\end{beamercolorbox}}}
